\section{零售明细的数据陷阱治理与假设边界}
\label{sec:data-governance}

\subsection{从原始字节到可计算业务实体}

原始文件不能依赖软件自动识别：其实际编码为GBK，订单日期与发货日期均按DD/MM/YYYY（Python格式\texttt{\%d/\%m/\%Y}）解释。若误用Latin-1，姓名和商品名会出现乱码；若沿用美式月日顺序，部分日期即使能够解析也会被静默交换。数据载入后统一转存为UTF-8，日期解析则固定格式，使后续月份归属与履约天数都建立在同一时间口径上。% src:交接/数据档案.json:文件编码,时间范围

业务实体审计比类型转换更关键。原表含$9800$条商品明细，却只对应$4921$个不同订单号，因此\texttt{Row ID}只作追溯键，订单规模按\texttt{Order ID}去重，常量\texttt{Country}不进入因素模型。% src:交接/数据档案.json:数据规模,列清单[名称=Order ID],数据陷阱
产品聚合固定使用\texttt{Product ID}：有$32$个产品编号对应多个名称，也有$16$个名称对应多个编号，冲突只标记而不凭名称合并。% src:求解/问题1/结果/数据治理审计.json:多名称ProductID数,多ProductID名称数
明细销售均值为$230.769$，中位数仅为$54.49$，故大额记录保留在原尺度总量中，典型水平与因素比较改用稳健统计或对数尺度。% src:交接/数据档案.json:数据陷阱[异常值]

排除唯一行号后，\texttt{Row ID}为$3406$与$3407$的记录其余字段完全相同。主口径删除$3407$并保留$3406$，得到$9799$条明细；敏感性口径保留原表，二者销售总额相差$281.372$。% src:求解/问题1/结果/数据治理审计.json:疑似复制Row ID,删除的Row ID,主口径行数,两口径总销售额差
这种处理只针对缺少数量字段佐证的唯一疑似复制组，后续总量结论仍用保留口径复核。

图\ref{fig:data-governance-flow}给出治理后的两条分析支路。明细解释层保留产品、客户、区域与发货模式，并以订单号形成重采样簇；月度预测层只按订单日期汇总销售额。发货日期异常另设双口径，不反向改变销售额或月份归属。% src:交接/图注素材.json:[图名=数据治理与口径分流图]

\begin{figure}[H]
  \centering
  \includegraphics[width=0.8\textwidth]{../求解/公共/图片/数据治理与口径分流图.png}
  \caption{零售明细治理与分析口径分流}
  \label{fig:data-governance-flow}
\end{figure}

\subsection{物流日期的双口径隔离}

发货日期早于下单日期的异常集中在$42$条明细、$20$个订单，均来自年末下单而发货年份写成早年。% src:求解/问题1/结果/数据治理审计.json:异常发货明细数,异常发货订单数
我们最初考虑把这些发货年份统一修正后作为唯一口径，但发现改为后一年后履约天数虽全部落在$3$至$7$天，年份仍未经企业确认，因此改用“年份修正”和“异常置缺”并行的双口径。% src:交接/实验记录.json:[问题=问题1/发货年份异常].尝试,现象,决定
设第$i$条明细的下单日为$d_i^O$，修正发货日与置缺发货日分别为$d_i^{S,c}$、$d_i^{S,m}$，履约天数定义为
\begin{equation}
L_i^{c}=d_i^{S,c}-d_i^O,\qquad
L_i^{m}=d_i^{S,m}-d_i^O.
\label{eq:dual-logistics-date}
\end{equation}
其中，$L_i^{c}$利用年份修正后的完整记录，$L_i^{m}$在异常行取缺失。两种处理均不修改\texttt{Order Date}、\texttt{Sales}、\texttt{Order ID}与\texttt{Ship Mode}，实际校验得到总销售额差、总订单数差及各发货模式销售额差均为零。% src:求解/问题1/结果/物流日期双口径.json:两口径总销售额差,两口径总订单数差,两口径发货模式销售额差
因此物流口径只约束履约关联的可信度，销售画像、驱动筛选的非物流字段与月度预测不会随修正选择而变化。

\subsection{贯穿四问的假设与影响边界}

表\ref{tab:model-assumption-boundary}把各问共用的假设压缩为可审计边界。每条假设都对应数据结构或字段缺口；一旦边界被突破，受影响的结论应回退到表中所列范围，而不是把局部口径扩展成普遍事实。

\begingroup
\footnotesize
\setlength{\tabcolsep}{3pt}
\begin{longtable}{>{\centering\arraybackslash}p{0.20\textwidth}>{\centering\arraybackslash}p{0.34\textwidth}>{\centering\arraybackslash}p{0.38\textwidth}}
\caption{建模假设、数据依据与影响边界}\label{tab:model-assumption-boundary}\\
\toprule
假设 & 本题依据 & 影响边界 \\
\midrule
\endfirsthead
\multicolumn{3}{c}{续表~\thetable}\\
\toprule
假设 & 本题依据 & 影响边界 \\
\midrule
\endhead
\bottomrule
\endfoot
销售可加，订单为最小业务簇 & 一单可含多条商品明细 & 总量按明细求和，重采样与切分按订单成簇 \\ % src:交接/数据档案.json:数据规模,列清单[名称=Order ID]
无记录月表示该层级无销售 & 研究期按订单日期形成连续月度记录 & 只适用于当前数据被视为全量明细；抽样数据须重估零填充波动 \\
治理判断保留敏感性口径 & 疑似复制缺少数量佐证，异常年份缺少企业确认 & 重复处理须双口径复核；日期修正只影响物流统计 \\ % src:求解/问题1/结果/数据治理审计.json:指标含义 % src:求解/问题1/结果/物流日期双口径.json:指标含义
历史顺序可用于外推 & 计划采用早期训练、后期封存检验 & 编码、参数与阈值只能由预测时点之前的数据估计 \\ % src:交接/计划.json:问题清单[编号=2].验证方案,问题清单[编号=3].验证方案
2026季节位置与份额机制延续，结构突变时触发滚动重估 & 季节总量以历史同月为低方差锚，层级份额由近期、同月与长期结构收缩得到 & 仅支持2026年既有类别与相近经营环境的外推；若季节次序、类别集合或份额路径突变，须滚动重估总量模型与收缩权重 \\ % src:交接/建模笔记_问题3.md:建模假设[弱结构延续假设];交接/建模笔记_问题4.md:建模假设[假设2,假设3];交接/结果解读_问题3.md:局限与使用边界
历史类别构成分解集合 & 层级预测在既有业务类别内分配总量 & 新类别须增设“其他/新层级”并重新归一化 \\
驱动只作关联解释 & 数据缺少数量、折扣、利润、成本、促销和库存 & 不推出利润最优、价格弹性、安全库存件数或标签干预收益 \\ % src:交接/数据档案.json:列清单,给建模师的话
\end{longtable}
\endgroup

这些边界决定后文的用词强度：产品与区域结果用于描述销售结构和预测信息，客户与物流层级可用于容量配置；任何经营建议都须经过时间外验证和敏感性检查，不能由单次分类均值直接推出干预收益。
